\documentclass[11pt]{article}
\usepackage[T1]{fontenc}
\usepackage{lmodern,amsmath,amssymb,amsthm}
\usepackage[margin=1in]{geometry}
\usepackage[colorlinks=true,linkcolor=blue,citecolor=blue,urlcolor=blue]{hyperref}
\newtheorem{theorem}{Theorem}
\newtheorem{proposition}[theorem]{Proposition}
\newtheorem{lemma}[theorem]{Lemma}
\newcommand{\BM}{d_{\mathrm{BM}}}
\newcommand{\Cb}{C_b}
\newcommand{\Hb}{H_b}
\title{Ordinary Banach--Mazur rigidity from bounded Hochschild vanishing}
\author{Alec Kriebel\\\small ORCID: \href{https://orcid.org/0009-0001-9320-500X}{0009-0001-9320-500X}}
\date{October 6, 2026\\\small Research draft: journal-source audit outstanding}
\hypersetup{pdftitle={Ordinary Banach-Mazur rigidity from bounded Hochschild vanishing},pdfauthor={Alec Kriebel},pdfsubject={Conditional consequence; dependency audit outstanding}}
\begin{document}
\maketitle
\begin{abstract}
Roydor's cohomological rigidity theorem and the vanishing assertion in
OpenAI's family 295 manuscript imply local ordinary complex
Banach--Mazur rigidity for every von Neumann algebra and its canonical
predual. We give the precise composition, including an elementary
adjoint argument for the predual conclusion. The threshold may depend on
the algebra and the conclusion is Jordan $*$-isomorphism. This draft
records a conditional consequence, rather than a certified new
unconditional theorem: Roydor's actual author slides have been read,
but his journal full text has not been obtained. The upstream proof
has received independent mathematical audits; a fresh kernel build
of its actual Lean declarations has not been reproduced. This
document is not cleared for Zenodo publication as a full solution.
\end{abstract}

\section{The target and conventions}

All Banach spaces and linear maps are complex. For Banach spaces $X,Y$, put
\[
 \BM(X,Y)=\inf_T \|T\|\,\|T^{-1}\|,
\]
where the infimum ranges over bounded bijective complex-linear maps
$T:X\to Y$ with bounded inverse. It is infinity when no such map
exists. The literal norm-product formula gives $\BM(0,0)=0$, and
$\BM(0,X)=\BM(X,0)=\infty$ for
$X\ne0$. For a von Neumann algebra $M$, $M_*$ is its canonical
Banach predual, with $M=(M_*)^*$ isometrically. The notation does
not assume a trace or a semifinite realization.

A Jordan $*$-isomorphism is a bijective complex-linear map preserving
involution and $x\circ y=(xy+yx)/2$. On nonzero unital algebras it
preserves the unit. The core question is whether, for each complex
von Neumann algebra $M$, there exists $\varepsilon_M>0$ such that
for every complex von Neumann algebra $N$ either inequality
\begin{equation}\label{eq:targets}
 \BM(M,N)<1+\varepsilon_M,
 \qquad \BM(M_*,N_*)<1+\varepsilon_M
\end{equation}
forces a Jordan $*$-isomorphism. The comparisons are with von Neumann
algebras and their canonical preduals. No assertion about arbitrary
nearby Banach spaces is intended.

\section{The two inputs}

For $k\ge1$, $\Cb^k(M,M)$ denotes bounded complex $k$-linear maps
with the usual multilinear operator norm; $\Cb^0(M,M)=M$. For a
$k$-cochain $f$, $k\ge1$, the differential is
\begin{align*}
 (df)(a_1,\ldots,a_{k+1})={}&a_1f(a_2,\ldots,a_{k+1})\\
 &+\sum_{j=1}^k(-1)^j
 f(a_1,\ldots,a_ja_{j+1},\ldots,a_{k+1})\\
 &+(-1)^{k+1}f(a_1,\ldots,a_k)a_{k+1}.
\end{align*}
In degree zero, $(db)(a)=ab-ba$. We use
\[
 \Hb^k(M,M)=\ker(d:\Cb^k\to\Cb^{k+1})/
 \operatorname{im}(d:\Cb^{k-1}\to\Cb^k).
\]
The denominator is the actual image, with no closure. Neither
complete boundedness nor normality is imposed on these cochains.

\begin{theorem}[Roydor's reduction, author-slide formulation]\label{thm:roydor}
Let $M$ be a complex von Neumann algebra with
$\Hb^2(M,M)=\Hb^3(M,M)=0$. There is $\varepsilon_M>0$ such that,
for every complex von Neumann algebra $N$, the following conditions
are equivalent: $M,N$ are Jordan $*$-isomorphic; $M,N$ are linearly
isometric; $M_*,N_*$ are linearly isometric; and each of the two
strict distance inequalities in \eqref{eq:targets} holds.
\end{theorem}

This is the stated scope of Theorem 2 in Roydor's CIRM slides,
interpreted in the standard complex operator-algebra setting,
\cite[PDF pp.~28--29]{RoydorSlides}, repeated later in the file. No
separability, factoriality, type or trace hypothesis occurs in that
statement. The bounded-cochain convention is introduced on PDF
pp.~30--32 and distinguished from completely bounded cohomology.
The displayed differential on p.~32 has apparent index/sign errors;
we use the standard differential written above, also used by family
295. The published article is \cite{Roydor2022}; its precise theorem
number, proof and correspondence with the slide version remain
unverified because journal full-text access was unsuccessful. We
therefore retain Theorem~\ref{thm:roydor} as an explicitly identified
external input rather than claiming to have independently reproved it.
The slides do not explicitly declare the scalar field; the journal's
preliminary field convention remains to be checked.

\begin{theorem}[OpenAI family 295 assertion]\label{thm:input}
For every complex von Neumann algebra $M$ and every $k\ge2$, every
bounded complex $k$-linear Hochschild cocycle $M^k\to M$ has a
bounded complex $(k-1)$-linear primitive. Thus $\Hb^k(M,M)=0$.
\end{theorem}

This is Theorem 1.1 of the manuscript \cite{OAI295}, pinned to repository
commit \texttt{adc7f124} (the full identifier is recorded in the
source receipt and the pinned link below). In particular,
its assertion has exactly the field, coefficient module, boundedness,
actual-image and degree conventions required above. Its proof reduces
ordinary cocycles to separately normal cocycles, constructs uniformly
bounded primitives on tracial type $\mathrm{II}_1$ algebras using a
unitary random walk and a Liouville argument, removes separability by
conditional expectations and compactness, and assembles central
summands with an explicit homotopy. The accompanying audit reports
record which arguments were checked. This description is an audit
summary, not a replacement proof of that input.

\section{The consequence and an independent predual argument}

\begin{lemma}\label{lem:adjoint}
For complex von Neumann algebras $M,N$,
\[
 \BM(M,N)\le\BM(M_*,N_*).
\]
\end{lemma}
\begin{proof}
If $T:M_*\to N_*$ is a bounded isomorphism, its Banach adjoint
$T^*:N\to M$ is a bounded complex-linear isomorphism, with inverse
$(T^{-1})^*$ and norms $\|T^*\|=\|T\|$,
$\|(T^*)^{-1}\|=\|T^{-1}\|$. Inverting this algebra-space
isomorphism does not change its distortion. Taking the infimum over
$T$ proves the inequality. If no $T$ exists, the right side is
infinite. Our zero convention gives the remaining cases.
\end{proof}

\begin{proposition}[Conditional composition]\label{prop:consequence}
If Theorems~\ref{thm:roydor} and \ref{thm:input} hold in the stated
scope, both targets in \eqref{eq:targets} hold for every complex von
Neumann algebra, with a common algebra-dependent threshold. In that
case Jordan $*$-isomorphism, algebra-space linear isometry, predual
linear isometry, and the two strict distance inequalities are equivalent.
\end{proposition}
\begin{proof}
For $M\ne0$, apply Theorem~\ref{thm:input} in degrees $2$ and $3$,
then apply Theorem~\ref{thm:roydor}. Its algebra threshold already
suffices for the predual implication by Lemma~\ref{lem:adjoint}.
Thus the predual conclusion needs no quantitative estimate beyond the
algebra reduction. The equivalences are supplied by the stated Roydor
input. For $M=0$, either strict distance inequality forces $N=0$,
so any positive threshold works.
\end{proof}

There is no use of an infimum attaining map. If an ordinary distance
is strictly below a given bound, some isomorphism has distortion
strictly below that bound, directly from the definition of infimum.
No continuity of the initial algebra-space isomorphism in a weak
topology is assumed. The predual map's adjoint is weak-star
continuous, but Roydor's ordinary-distance assertion does not require
that extra property.

\section{Scope, attribution and unresolved verification}

The outcome sought is Jordan $*$-isomorphism. Opposite algebras have
the same underlying Banach norm and are Jordan $*$-isomorphic, so
ordinary Banach-space geometry cannot determine an orientation of the
associative product. Likewise, no universal $\varepsilon$ follows from
the qualitative input: its quantifiers are $\forall M\ \exists
\varepsilon_M>0\ \forall N$, not a uniform threshold over $M$.

Ricard and Roydor already proved unconditional stability for
\emph{completely bounded} Banach--Mazur distance
\cite[Theorem A]{RicardRoydor}. Completely bounded distortion involves
all matrix levels; its smallness is a stronger hypothesis than the
ordinary distortion considered here. The earlier theorem therefore
does not, by itself, supply the ordinary-distance conclusion.
This draft does not claim a new completely bounded result.

Roydor's slides already state ordinary rigidity when the type
$\mathrm{II}_1$ summand has property $\Gamma$, a Cartan subalgebra,
or McDuff absorption \cite[PDF pp.~47--48]{RoydorSlides}. Those
subcases are inherited. Ricard--Roydor also record a conditional
two-matrix-level distance version \cite[Remark 3.7]{RicardRoydor};
that stronger metric hypothesis is distinct from ordinary distance.

The claimed cohomology vanishing is attributed to OpenAI, and the
cohomological geometric reduction to Roydor. The present composition
is an immediate consequence if both dependencies are correct, not
an independent solution of the Kadison--Ringrose cohomology problem
or a new rigidity mechanism. The adjoint observation is standard
Banach-space duality. A bounded current primary-source search and
search of the pinned upstream manuscript corpus found no duplicate
general ordinary consequence. Absence of a search result does not
establish priority or firstness. The original reduction was publicly
described before its journal publication; a 2020 slide date is not
evidence for the upstream vanishing assertion's public priority.

The present package preserves a mathematical proof audit of family
295 and a static inspection of its genuine Lean declarations. It
does not claim that a formalization catalogue or a comparator stub
verifies a theorem. A fresh Lean kernel reproduction remains pending,
and the follow-on Banach--Mazur theorem is not formalized here.
The independent mathematical source audit is the intended validation
basis for the cohomology input; a missing kernel build is a limitation
of formal evidence, not a mathematical counterexample.
Full journal-source validation of Roydor's reduction, completion of
the priority audit, and fresh reviews of a final publication package
are still required. No Zenodo deposit or tracker entry is warranted
as a full resolution at this stage.

\paragraph{Disclosure.}
AI tools were used extensively for research, drafting and
verification. Independent automated adversarial reviews are internal
checks, not conventional human peer review. This research draft has
not undergone conventional human peer review or refereeing.

\begin{thebibliography}{9}
\bibitem{Roydor2022}
J.~Roydor, \emph{Banach--Mazur stability of von Neumann algebras},
Journal of Topology and Analysis \textbf{14}(3) (2022), 767--792.
\href{https://doi.org/10.1142/S1793525321500151}{doi:10.1142/S1793525321500151}.
Published online December 11, 2020. Journal full text not read in this audit.
\bibitem{RoydorSlides}
J.~Roydor, \emph{Banach--Mazur stability of von Neumann algebras},
CIRM lecture slides, October 2020, Theorem 2, PDF pp.~28--29.
\href{https://www.cirm-math.fr/RepOrga/2169/Slides/Roydor_Slides.pdf}{Primary author slides}.
\bibitem{OAI295}
OpenAI, \emph{Vanishing of higher bounded Hochschild cohomology},
OpenAI Math Release preprint, manuscript dated September 23, 2026.
\href{https://github.com/openai/math/blob/adc7f1241b42e322a6451854ab7e4b4c146bf78a/preprints/Vanishing-of-higher-bounded-Hochschild-cohomology-September-23-2026/paper.pdf}{Pinned source and paper}.
Citation metadata supplied by the manuscript-specific README.
\bibitem{RicardRoydor}
\'{E}.~Ricard and J.~Roydor, \emph{A non-commutative Amir--Cambern
theorem for von Neumann algebras and nuclear $C^*$-algebras},
Journal of Functional Analysis \textbf{267}(4) (2014), 1121--1136,
\href{https://doi.org/10.1016/j.jfa.2014.05.018}{doi:10.1016/j.jfa.2014.05.018};
\href{https://arxiv.org/abs/1108.1970v2}{arXiv:1108.1970v2}.
\end{thebibliography}
\end{document}
